% Publication insert.  The main manuscript inputs this file once; it in turn
% inputs the repaired CMP(1) and internal-sheet closure files.

\subsection{Thermodynamic carrier reduction at fixed lattice spacing}
\label{sec:g18-thermodynamic-reduction}

The finite-volume Riesz island first has a thermodynamic consequence that is
strictly weaker than a particle theorem but stronger than a finite-volume rank
statement.  We isolate that projection step before proving the stronger band,
source-frame and patchwise carrier theorems below; no survival of a
finite-volume eigenvalue is silently assumed.

Let $\mathcal A_{\mathrm{loc}}$ be the local observable algebra of the
unconstrained link system on $\mathbb Z^3$.  Write
$(\mathcal H_u,\pi_u,\Omega_u)$ for the GNS triple of its interacting
ground state, and let $H_\infty(u)$ be the GNS Hamiltonian normalized by
$H_\infty(u)\Omega_u=0$.  Put
\[
 E_F=\frac83,\qquad
 I_F(u)=\left[\frac83(1-\lambda),\frac83(1+\lambda)\right],
 \qquad \lambda=27c_2|u| .
\]

\begin{proposition}[Infinite-volume external Riesz projection]
\label{prop:g18-infinite-riesz}
Assume
\[
 27|u|<c_1,\qquad \lambda<\frac1{33}.
\]
Then Theorems~2 and~3 of~\cite{Yarotsky2004QP}, applied to the same
three-link cell decomposition as Proposition~\ref{prop:periodic-yarotsky-descent},
give a translation-invariant thermodynamic ground state and a self-adjoint
$H_\infty(u)$ satisfying
\begin{equation}
 \operatorname{spec}H_\infty(u)
 \subset
 \bigcup_{a\in\operatorname{spec}H_{E,\infty}^{\rm kin}}
 [a(1-\lambda),a(1+\lambda)] .
 \label{eq:g18-infinite-enclosure}
\end{equation}
Consequently the contour $\Gamma_F$ in~\eqref{eq:riesz-contour} lies in the
resolvent set of $H_\infty(u)$ and defines
\begin{equation}
 P_{F,\infty}(u)
 =-\frac1{2\pi i}\oint_{\Gamma_F}
       (H_\infty(u)-z)^{-1}\,dz .
 \label{eq:g18-infinite-projection}
\end{equation}
The target interval is externally separated from the rest of the spectrum
by
\[
 \delta_-(u)\ge\frac{1-15\lambda}{3},\qquad
 \delta_+(u)\ge\frac{1-33\lambda}{6}.
\]
For local $A,B\in\mathcal A_{\mathrm{loc}}$, along Yarotsky's open-box
thermodynamic sequence (or a symmetry-preserving boundary sequence for which
the same common enclosure and weak-resolvent limit have separately been
proved),
\begin{equation}
 \lim_{L\to\infty}
 \langle A\Omega_L,P_{F,L}(u)B\Omega_L\rangle
 =
 \langle \pi_u(A)\Omega_u,
 P_{F,\infty}(u)\pi_u(B)\Omega_u\rangle .
 \label{eq:g18-projection-limit}
\end{equation}
The local gauge-invariant, charge-odd subspace is reducing, so all of these
statements remain true after restriction to the physical charge-odd local
GNS sector.
\end{proposition}

\begin{proof}
The rescaling and local norm estimate in
Proposition~\ref{prop:periodic-yarotsky-descent} verify the hypotheses of
\cite[Theorems~2 and~3]{Yarotsky2004QP}.  Those theorems give the weak-star
ground-state limit, the weak-resolvent limit, and the same multiplicative
spectral enclosure with the finite-volume free spectrum replaced by its
finite-support additive semigroup.  The Casimir arithmetic in the proof of
Theorem~\ref{thm:riesz-island} is unchanged for that semigroup: the nearest
free energies to $16/6$ are $14/6$ and $17/6$.  This proves the displayed
gaps and places $\Gamma_F$ in a common resolvent set.

Yarotsky's weak-resolvent convergence is initially stated for
$z\in\mathbb C\setminus\mathbb R$.  Choose one connected tubular
neighbourhood of $\Gamma_F$ disjoint from every finite- and infinite-volume
spectral enclosure.  Matrix elements of the resolvents are uniformly bounded
there by the common spectral clearance.  The normal-family theorem and the
identity theorem therefore extend their convergence to the contour.
Integrating the matrix elements around $\Gamma_F$ proves
~\eqref{eq:g18-projection-limit}.

Invariance of the thermodynamic state under finite-support gauge
transformations and charge conjugation gives their canonical GNS
implementing unitaries.  The corresponding automorphisms commute with the
dynamics, so the implementers commute with $H_\infty(u)$.  Define the physical
odd sector to be the closed cyclic subspace generated by local
gauge-invariant, charge-odd observables.  It is reducing for the Hamiltonian,
its resolvent and the contour projection, which proves the last assertion.
\end{proof}

\begin{remark}[What the enclosure alone does not say]
\label{rem:g18-projection-may-vanish}
The enclosure is one-sided: it says where spectrum may occur, not that every
free patch survives.  Thus Proposition~\ref{prop:g18-infinite-riesz} by
itself does not prove $P_{F,\infty}(u)\ne0$.  Nor does the finite-volume rank
$3L^3$ imply a nonzero matrix element against any fixed local source.  The
next two results supply the missing close-projection, totality and source
estimates; they are not consequences of the enclosure by itself.
\end{remark}

% Publication-ready replacement for the unsupported CMP(1) paragraph in
% G18_FIXED_SPACING_CARRIER_BRIDGE_INSERT.tex, lines 103--184.
% This file is input by G18_FIXED_SPACING_CARRIER_BRIDGE_INSERT.tex.

\subsubsection{Coefficient close projection and transport to the physical GNS sector}
\label{sec:g18-cmp1-complete}

We use the three-outgoing-link cell decomposition and the gap-one
normalization of Proposition~\ref{prop:periodic-yarotsky-descent}.  If
$J\Subset\mathbb Z^3$ is a nonempty exact cell support, write
$\mathcal H'_J$ for the corresponding excited tensor factor and
$H_{J,0}$ for the restriction of the free cell Hamiltonian.  Thus
\begin{equation}
 H_{J,0}\ge |J|\,1_{\mathcal H'_J}.
 \label{eq:g18-cmp1-cell-gap}
\end{equation}
The superscript minus below means both local gauge invariance and odd charge
conjugation parity.  Define
\begin{align}
 \mathcal X_0^-&=
 \left\{v=(v_J):v_J\in\mathcal H'_J,
       \ \|v\|_0:=\sum_J\|v_J\|<\infty\right\}^-,
 \label{eq:g18-cmp1-x0}\\
 \mathcal X_1^-&=
 \left\{v\in\mathcal X_0^-:v_J\in\operatorname{Dom}H_{J,0},
       \ \|v\|_1:=\sum_J\|H_{J,0}v_J\|<\infty\right\}.
 \label{eq:g18-cmp1-x1}
\end{align}
The empty coefficient is absent in the odd sector.  Consequently
\eqref{eq:g18-cmp1-cell-gap} makes $\|\cdot\|_1$ a norm, and the finite
coefficient families with entries in $\operatorname{Dom}H_{J,0}$ form a
core, denoted $\mathcal c_{00}^-$, for
\begin{equation}
 D(v_J)_J=(H_{J,0}v_J)_J,
 \qquad \operatorname{Dom}D=\mathcal X_1^-.
 \label{eq:g18-cmp1-D}
\end{equation}

\begin{lemma}[Exact-support symmetry]
\label{lem:g18-cmp1-exact-support-symmetry}
The exact-support projections commute with every local gauge
transformation and with charge conjugation.  Hence $D$, Yarotsky's
coefficient perturbation $F_u$, the free shell projection $R^-$, and the
coefficient Riesz projection $\mathcal P_u^-$ preserve
$\mathcal X_j^-$, $j=0,1$.
\end{lemma}

\begin{proof}
Let $p_x=|\Omega_x\rangle\langle\Omega_x|$ and $q_x=1-p_x$ on a grouped
cell.  In a finite box the exact-support projector is a tensor product of
$q_x$ for $x\in J$ and $p_x$ for $x\notin J$.  A gauge transformation at a
vertex acts on incident link factors by left or right regular
representations.  After the links are regrouped into cells this action still
factorizes over the cell tensor factors, and every factor fixes the constant
cell vacuum $\Omega_x$.  It therefore commutes with $p_x$, $q_x$, and every
exact-support projector.  Charge conjugation is the linkwise pullback by the
group automorphism $U\mapsto\overline U$; it also factorizes and fixes
$\Omega_x$, so the same conclusion holds.

 The free electric Hamiltonian and the plaquette interaction commute with
 these symmetries.  The symmetries fix the product vacuum and preserve the
 normalization
 $\langle\widetilde\Omega_\Lambda(u),\Omega_{\Lambda,0}\rangle=1$.
 Since the interacting ground line is unique in the small-coupling domain,
 this normalization fixes its representative and makes that representative
 invariant.  Uniqueness of the finite-volume coefficient expansion then
 implies that $D$ and $F_u$
commute with the induced coefficient actions.  Passing to Yarotsky's
thermodynamic coefficient limits preserves these commutators.  Functional
calculus for $D$ and the contour definition of $\mathcal P_u^-$ prove the
last assertion.
\end{proof}

\begin{lemma}[Unweighted coefficient close projection]
\label{lem:g18-cmp1-close-projection-complete}
Let $g=\sup_x\|\widehat\phi_x(u)\|=27|u|$.  There is a function
$a(g)\to0$ such that
\begin{equation}
 \|F_uv\|_0\le a(g)\|v\|_1,
 \qquad v\in\mathcal X_1^-.
 \label{eq:g18-cmp1-relative-bound}
\end{equation}
 Assume in addition the spectral-enclosure conditions
 \begin{equation}
  27|u|<c_1,\qquad
  \lambda:=27c_2|u|<\frac1{33}.
  \label{eq:g18-cmp1-enclosure-smallness}
 \end{equation}
 On the circle
 \begin{equation}
  \gamma_u=\{z:|z-4|=r_u\},
  \qquad r_u=\frac{1-\lambda}{8},
  \label{eq:g18-cmp1-dynamic-contour}
 \end{equation}
 put
 \begin{equation}
  R^-=-\frac1{2\pi i}\oint_{\gamma_u}(D-z)^{-1}\,dz,
  \qquad
  \mathcal P_u^-=-\frac1{2\pi i}
        \oint_{\gamma_u}(D+F_u-z)^{-1}\,dz .
  \label{eq:g18-cmp1-projections}
 \end{equation}
 If $q(g):=35a(g)<1$, then
\begin{equation}
 \boxed{\quad
 \|\mathcal P_u^--R^-\|_{\mathcal X_0^-\to\mathcal X_0^-}
 \le \kappa(g):=\frac{q(g)}{1-q(g)}.
 \quad}
 \label{eq:g18-cmp1-kappa}
\end{equation}
In particular, $\kappa(g)\to0$; after decreasing the small-coupling
interval, $\kappa(g)<1$.
\end{lemma}

\begin{proof}
Fix any $0<\rho<1$.  Yarotsky's estimate~(14), for an input of exact
support $J$, gives
\[
 \sum_K\|(F_uv_J)_K\|\rho^{-(d(K,J)+1)}
 \le C_{14}(\rho)g|J|\|v_J\|.
\]
Since the weight on the left is at least $\rho^{-1}$,
\[
 \sum_K\|(F_uv_J)_K\|
 \le \rho C_{14}(\rho)g|J|\|v_J\|
 \le \rho C_{14}(\rho)g\|H_{J,0}v_J\|.
\]
Summing the columns proves~\eqref{eq:g18-cmp1-relative-bound} with
$a(g)=\rho C_{14}(\rho)g$.  This also proves that $D+F_u$, with domain
$\mathcal X_1^-$, is closed when $a(g)<1$.

 The retained-shell gaps put $\gamma_u$ in the resolvent set of $D$ and give
 \[
  K_0:=\sup_{z\in\gamma_u}\|(D-z)^{-1}\|_{0\to0}
  \le\frac8{1-\lambda}.
 \]
 Moreover,
 \[
  D(D-z)^{-1}=1+z(D-z)^{-1},
  \qquad
  K_1:=\sup_{z\in\gamma_u}\|D(D-z)^{-1}\|_{0\to0}
  \le2+\frac{32}{1-\lambda}<35,
 \]
 because $|z|\le4+r_u$ on $\gamma_u$.  Therefore
\[
 \|F_u(D-z)^{-1}\|_{0\to0}\le a(g)K_1\le q(g).
\]
For $q(g)<1$, Yarotsky's Neumann series~(15) converges in
$\mathcal X_0^-$ and yields
\[
 \|(D+F_u-z)^{-1}-(D-z)^{-1}\|_{0\to0}
 \le K_0\frac{q(g)}{1-q(g)}.
\]
 The length of $\gamma_u$ is $2\pi r_u$, and
 $r_uK_0\le1$.  Contour integration therefore proves
 \eqref{eq:g18-cmp1-kappa}.
\end{proof}

\begin{lemma}[Finite-support shell bridge]
\label{lem:g18-cmp1-finite-support-shell-bridge}
In the physical charge-odd coefficient space,
\begin{equation}
 \operatorname{spec}(D|_{\mathcal X_0^-})\cap(-\infty,5)=\{4\},
 \qquad
 \operatorname{Ran}R^-
 =\overline{\operatorname{span}}^{\,\|\cdot\|_0}
   \{w_{x,\alpha}:x\in\mathbb Z^3,\ \alpha=1,2,3\}.
 \label{eq:g18-cmp1-free-shell-span}
\end{equation}
\end{lemma}

\begin{proof}
Let $v_J$ be a finite-support gauge-invariant, charge-odd eigen-coefficient
of rescaled electric energy below $5$.  Choose a periodic side length $L$
larger than the occupied-edge bounding box and embed the support injectively
in the $L^3$ torus, extending by the cell vacuum.  No occupied path winds
  around the torus, so the embedded coefficient belongs to the trivial
  electric one-form sector.  The retained-shell theorem then forces its energy
  to be $4$ and its energy-$4$ component to be a linear combination of odd
  elementary plaquette vectors.  The cell Laplacians have discrete Casimir
  spectrum, and energy below $5$ bounds the number of occupied cells.  Thus
  this diagonal direct sum has no additional approximate spectrum below $5$.
  Finite coefficient families are a core for
$D$, and spectral truncation in each finite tensor factor followed by
truncation in $J$ is dense in $\mathcal X_0^-$.  Applying these truncations
to the spectral projection of $D$ proves
\eqref{eq:g18-cmp1-free-shell-span}.
\end{proof}

\begin{lemma}[Coefficient totality]
\label{lem:g18-cmp1-coefficient-totality-complete}
Assume $\kappa(g)<1$.  Then
\begin{equation}
 \operatorname{Ran}\mathcal P_u^-
 =\overline{\mathcal P_u^-\operatorname{Ran}R^-}^{\,\|\cdot\|_0}.
 \label{eq:g18-cmp1-two-projection-totality}
\end{equation}
Using the finite-support retained-shell classification,
\begin{equation}
 \operatorname{Ran}\mathcal P_u^-
 =\overline{\operatorname{span}}^{\,\|\cdot\|_0}
 \{\mathcal P_u^-w_{x,\alpha}:
 x\in\mathbb Z^3,\ \alpha=1,2,3\}.
 \label{eq:g18-cmp1-seed-totality}
\end{equation}
\end{lemma}

\begin{proof}
Write $P=\mathcal P_u^-$, $R=R^-$, and $Q=1-R$.  Both $R$ and $Q$ are
blockwise orthogonal projections and hence contractions in the direct-sum
$\ell^1$ norm.  For $v\in\mathcal X_0^-$ set
\[
 y_0=Qv,
 \qquad y_{n+1}=QPy_n.
\]
Every $y_n$ lies in $\operatorname{Ran}Q$, and hence
\begin{equation}
 \|y_{n+1}\|_0
 =\|Q(P-R)y_n\|_0
 \le\kappa(g)\|y_n\|_0.
 \label{eq:g18-cmp1-geometric-irrelevant}
\end{equation}
Idempotence of $P$ gives the exact identity
\begin{equation}
 Py_n=P^2y_n=PRPy_n+PQPy_n=PRPy_n+Py_{n+1}.
 \label{eq:g18-cmp1-iteration-identity}
\end{equation}
Consequently, for every $N\ge0$,
\begin{equation}
 Pv=PRv+\sum_{n=0}^{N}PRPy_n+Py_{N+1}.
 \label{eq:g18-cmp1-iteration-finite}
\end{equation}
The last term tends to zero by
\eqref{eq:g18-cmp1-geometric-irrelevant}.  Every preceding term belongs to
$P\operatorname{Ran}R$, which proves~\eqref{eq:g18-cmp1-two-projection-totality}.
The reverse inclusion is immediate because $P$ is a projection.

 Lemma~\ref{lem:g18-cmp1-finite-support-shell-bridge} says that
 $\operatorname{Ran}R^-$ is the $\ell^1$-closed span of the
 $w_{x,\alpha}$.
Applying the bounded operator $P$ and using
\eqref{eq:g18-cmp1-two-projection-totality} proves
\eqref{eq:g18-cmp1-seed-totality}.  This is the complete version of the
iteration compressed into equations~(21)--(24) of
\cite{Yarotsky2004QP}; no one-site nondegeneracy assumption is used.
\end{proof}

\begin{lemma}[Coefficient-to-GNS transport]
\label{lem:g18-cmp1-gns-transport-complete}
 Let $\mathcal H_{u,\mathrm{phys}}^-$ be the closed local gauge-invariant,
 charge-odd cyclic subspace in the GNS representation of the thermodynamic
 ground state.  Put
 \begin{equation}
  \widehat H_\infty^-(u)=\frac32H_\infty^-(u),
  \label{eq:g18-cmp1-rescaled-H-infty}
 \end{equation}
 and let $P_{F,\infty}^-(u)$ be its shell Riesz projection on
 $\gamma_u$.  This is the same projection as the shell Riesz projection of
 $H_\infty^-(u)$, because multiplication of the Hamiltonian and contour by
 the same positive scalar does not change a spectral projection.
Then the coefficient map extends to a contraction
\begin{equation}
 \Gamma_u:\mathcal X_0^-\longrightarrow
 \mathcal H_{u,\mathrm{phys}}^-,
 \qquad
 \Gamma_uv=\pi_u\!\left(\sum_J\widehat v_J\right)\Omega_u,
 \label{eq:g18-cmp1-Gamma}
\end{equation}
and satisfies
\begin{equation}
 P_{F,\infty}^-(u)\Gamma_u
 =\Gamma_u\mathcal P_u^-.
 \label{eq:g18-cmp1-riesz-intertwining}
\end{equation}
Moreover,
\begin{equation}
 \operatorname{Ran}P_{F,\infty}^-(u)
 =\overline{\operatorname{span}}
 \{P_{F,\infty}^-(u)
       \pi_u(\widehat w_{x,\alpha})\Omega_u\}.
 \label{eq:g18-cmp1-gns-totality}
\end{equation}
No injectivity assumption on $\Gamma_u$ is required.
\end{lemma}

\begin{proof}
For a finite coefficient family,
\[
 \|\Gamma_uv\|
 \le\sum_J\|\pi_u(\widehat v_J)\Omega_u\|
 \le\sum_J\|\widehat v_J\|
 =\sum_J\|v_J\|=\|v\|_0.
\]
Thus the series in~\eqref{eq:g18-cmp1-Gamma} converges in operator norm and
$\Gamma_u$ extends by continuity to $\mathcal X_0^-$.  Exact-support
symmetry, Lemma~\ref{lem:g18-cmp1-exact-support-symmetry}, puts its range in
the physical odd sector.

 Yarotsky's equation~(18), in the present gap-one normalization, gives on the
 finite coefficient core
 \begin{equation}
  \widehat H_\infty^-(u)\Gamma_uv=\Gamma_u(D+F_u)v.
 \label{eq:g18-cmp1-generator-intertwining}
\end{equation}
Finite exact-support and local spectral truncations approximate every
$v\in\mathcal X_1^-$ in the $D$-graph norm.  The relative bound
\eqref{eq:g18-cmp1-relative-bound}, boundedness of $\Gamma_u$, and closedness
 of $\widehat H_\infty^-(u)$ therefore extend
\eqref{eq:g18-cmp1-generator-intertwining} to all of
 $\mathcal X_1^-$.  If $z$ lies on $\gamma_u$ and $a\in\mathcal X_0^-$,
apply this equality to
$v=(D+F_u-z)^{-1}a$ to obtain
\begin{equation}
  (\widehat H_\infty^-(u)-z)^{-1}\Gamma_ua
 =\Gamma_u(D+F_u-z)^{-1}a.
 \label{eq:g18-cmp1-resolvent-intertwining}
\end{equation}
Contour integration proves~\eqref{eq:g18-cmp1-riesz-intertwining}.

 It remains only to justify density in the correct sector.  The paragraph
 preceding equation~(18) of \cite{Yarotsky2004QP} states that finite linear
 combinations of creator vectors $\pi_u(\widehat v_J)\Omega_u$ are dense in
 the GNS Hilbert space.  To restrict this dense set to the physical odd
 sector, average a finite creator family over
the compact product of gauge groups at the finitely many vertices incident
to its support, and then apply $(1-C)/2$.  By
Lemma~\ref{lem:g18-cmp1-exact-support-symmetry}, this operation neither
enlarges exact support nor leaves the finite coefficient core.  Since the
corresponding Hilbert-space average is an orthogonal projection that fixes
every physical odd vector, the averaged creator families are dense in
$\mathcal H_{u,\mathrm{phys}}^-$.  Hence
\[
 \operatorname{Ran}P_{F,\infty}^-(u)
 =\overline{P_{F,\infty}^-(u)\Gamma_u\mathcal c_{00}^-}
 =\overline{\Gamma_u\mathcal P_u^-\mathcal c_{00}^-}.
\]
Now apply Lemma~\ref{lem:g18-cmp1-coefficient-totality-complete} and the
boundedness of $\Gamma_u$ to obtain
\eqref{eq:g18-cmp1-gns-totality}.
\end{proof}

\begin{corollary}[Uniform nonvanishing of each canonical seed]
\label{cor:g18-cmp1-seed-nonzero-complete}
For all sufficiently small $g$, every vector
\[
 \psi_{x,\alpha}
 =P_{F,\infty}^-(u)\pi_u(\widehat w_{x,\alpha})\Omega_u
\]
is nonzero, uniformly in $x$ and $\alpha$.
\end{corollary}

\begin{proof}
 The normalized creator has three-cell support $S_s$ and obeys
 $\widehat w_s^{\,*}\widehat w_s
 =p_{S_s}:=|\Omega_{S_s,0}\rangle\langle\Omega_{S_s,0}|\otimes I$.
 Here is the needed local product-vacuum estimate.  In a periodic volume,
 the product vacuum is a trial state, while the sum of interaction terms is
 bounded below by $-|\Lambda|g$.  Hence
 \[
  E_\Lambda(u)\le |\Lambda|g,
  \qquad
  E_\Lambda(u)\ge
  \sum_{x\in\Lambda}\langle\widehat h_x\rangle_{\Lambda,u}
       -|\Lambda|g .
 \]
 Translation invariance gives
 $\langle\widehat h_x\rangle_{\Lambda,u}\le2g$.  Since
 $1-p_x\le\widehat h_x$ and
 $1-p_{S_s}\le\sum_{x\in S_s}(1-p_x)$, every periodic subsequential limit
 obeys
 \[
  1-\omega_u(p_{S_s})\le2|S_s|g=6g.
 \]
 Yarotsky uniqueness identifies that limit with $\omega_u$.
Since $R^-w_s=w_s$, Lemma~\ref{lem:g18-cmp1-close-projection-complete} and
the contraction property of $\Gamma_u$ imply
\begin{align*}
 \|\psi_s\|
 &=\|\Gamma_u\mathcal P_u^-w_s\|\\
 &\ge \|\Gamma_uw_s\|
       -\|\Gamma_u(\mathcal P_u^--R^-)w_s\|\\
 &\ge \sqrt{1-6g}-\kappa(g).
\end{align*}
The last quantity is positive after decreasing the small-coupling interval.
\end{proof}

\begin{remark}[Scope]
\label{rem:g18-cmp1-scope-complete}
Lemmas~\ref{lem:g18-cmp1-close-projection-complete}--
\ref{lem:g18-cmp1-gns-transport-complete} close CMP(1) for the complete
three-component shell at fixed lattice spacing.  They do not identify an
individual internal dispersion sheet, an isotropic Wilson transfer-matrix
particle, or a continuum Yang--Mills particle.
\end{remark}


\begin{lemma}[Exponentially rooted shell dressing]
\label{lem:g18-rooted-shell-dressing}
Under the hypotheses of
Lemma~\ref{lem:g18-cmp1-close-projection-complete}, fix $0<\eta<1$ small
enough for Yarotsky's weighted estimate~(14).  For every seed
$s=(x,\alpha)$,
\begin{equation}
 \mathcal P_u^-w_s=w_s+c_s,\qquad
 \sum_J\eta^{-d(J,S_s)}
       \|H_{J,0}c_{s,J}\|
 \le A_\eta(g),
 \qquad A_\eta(g)\xrightarrow[g\to0]{}0 ,
 \label{eq:g18-rooted-projection}
\end{equation}
uniformly in $s$ and in finite-volume approximants.
\end{lemma}

\begin{proof}
For a root set $S$, put
\[
 \|v\|_{S,0}^{(\eta)}
 =\sum_J\eta^{-d(J,S)}\|v_J\|,
 \qquad
 \|v\|_{S,1}^{(\eta)}
 =\sum_J\eta^{-d(J,S)}\|H_{J,0}v_J\|.
\]
Yarotsky's estimate~(14), the triangle inequality
$d(K,S)\le d(K,J)+d(J,S)$, and $H_{J,0}\ge |J|$ give
\[
 \|F_uv\|_{S,0}^{(\eta)}
 \le a_\eta(g)\|v\|_{S,1}^{(\eta)},
 \qquad a_\eta(g)=\eta C_{14}(\eta)g.
\]
On the dynamic contour $\gamma_u$ of
\eqref{eq:g18-cmp1-dynamic-contour}, put
\[
 K_1=\sup_{z\in\gamma_u}
 \|D(D-z)^{-1}\|_{S,0\to S,0}<35,
 \qquad q_\eta=a_\eta(g)K_1.
\]
For $q_\eta<1$, the Neumann resolvent series converges in the rooted norm.
Since $R^-w_s=w_s$, applying $D$ to the resolvent difference and integrating
gives
\[
 A_\eta(g)\le r_uK_1\frac{q_\eta}{1-q_\eta}.
\]
This proves~\eqref{eq:g18-rooted-projection}.  The argument accepts an
arbitrary finite input support and does not use the one-site nondegeneracy
assumption of \cite[Theorem~4]{Yarotsky2004QP}.
\end{proof}

\begin{theorem}[Fixed-spacing three-component physical band and Wilson-source frame]
\label{thm:g18-fixed-spacing-band}
 Under the hypotheses of Proposition~\ref{prop:g18-infinite-riesz}, shrink
 the existential small-$|u|$ interval so that
 Lemmas~\ref{lem:g18-cmp1-close-projection-complete} and
 \ref{lem:g18-rooted-shell-dressing} hold.  In the physical charge-odd GNS
sector,
\begin{equation}
 \operatorname{Ran}P_{F,\infty}^-(u)
 \simeq \ell^2(\mathbb Z^3)\otimes\mathbb C^3
 \label{eq:g18-three-component-band}
\end{equation}
by a translation-intertwining unitary.  In this representation the
Hamiltonian is convolution by an exponentially summable Hermitian
$3\times3$ kernel, so its Bloch matrix is analytic in a complex momentum
strip.  The complete band retains the external gaps in
Proposition~\ref{prop:g18-infinite-riesz}.

More precisely, let
\[
 \psi_{x,\alpha}
 =P_{F,\infty}^-(u)\pi_u(\widehat w_{x,\alpha})\Omega_u,
 \qquad
 G_{st}=\langle\psi_s,\psi_t\rangle ,
\]
where $\widehat w_{x,\alpha}$ is the bounded rank-one plaquette creator.
For some $\mu>0$,
\begin{equation}
 \|G-I\|_{\mu,\sharp}=g_0(u)<1,\qquad
 G(k)\succeq(1-g_0(u))I_3
 \quad\hbox{for every real }k .
 \label{eq:g18-proved-canonical-frame}
\end{equation}

The same conclusion holds for the literal three-orientation Wilson
multiplication source
\[
 O^W_{x,\alpha}=\operatorname{ImTr}U_{p(x,\alpha)}.
\]
After the fixed normalization $c_W=i\sqrt2$, its projected synthesis map
$T_W$ is an isomorphism onto the band and its Gram symbol satisfies
\begin{equation}
 G_W(k)\succeq
 (1-\|G-I\|_{\mu,\sharp})
 (1-\|E\|_{\mu,\sharp})^2I_3>0
 \quad\hbox{for every real }k ,
 \label{eq:g18-wilson-frame}
\end{equation}
where $\|E\|_{\mu,\sharp}\to0$ as $u\to0$.
\end{theorem}

\begin{proof}
 Lemma~\ref{lem:g18-cmp1-gns-transport-complete} gives
 \[
 \operatorname{Ran}P_{F,\infty}^-
 =\overline{\operatorname{span}}\{\psi_{x,\alpha}\}.
\]
It remains to prove that these generators form a frame rather than merely a
dense family.

Yarotsky's cluster estimate~(16) has the form
\[
 |\omega_u(A^*B)-\omega_u(A^*)\omega_u(B)|
 \le C_0^{|I|+|J|}\tau(g)^{\operatorname{dist}(I,J)}
       \|A\|\|B\|,\qquad \tau(g)\to0 .
\]
The unit onsite gap and the variational energy-density bound also give, for
every fixed cell set $T$,
\[
 1-\omega_u(Q_T)\le2|T|g,\qquad
 |\omega_u(A)-\omega_0(A)|
 \le2\|A\|\sqrt{2|T|g}+4|T|g\|A\|.
\]
Use orthogonality of $P_{F,\infty}^-$ to expand only one Riesz mark:
\[
 G_{st}
 =\omega_u(\widehat w_s^*\widehat w_t)
  +\sum_J\omega_u(\widehat w_s^*\widehat c_{t,J}).
\]
Charge oddness kills all one-point functions and the empty coefficient.
For separated roots, the cluster estimate and
~\eqref{eq:g18-rooted-projection} give
\[
 |G_{st}-\delta_{st}|
 \le C\{\tau(g)^{R_{st}}+
 A_\eta(g)\theta^{R_{st}}\},
 \qquad \theta=\max\{C_0\eta,\tau(g)\}<1 .
\]
For the finitely many contacting roots, local vacuum continuity gives
$|G_{st}-\delta_{st}|\le C_{\rm loc}\sqrt g+A_\eta(g)$.
Choose $\eta$ first and then $g$ so that
$e^\mu C_0\eta<1$ and $e^\mu\tau(g)<1$.  The number of roots at distance
$n$ grows only quadratically, hence the exponentially weighted row and
column sums tend to zero with $g$.  This proves
~\eqref{eq:g18-proved-canonical-frame}.

The synthesis $J_P$ therefore obeys
\[
 (1-g_0)\|f\|_2^2\le\|J_Pf\|^2\le(1+g_0)\|f\|_2^2.
\]
Its range is closed and, by totality, is the whole Riesz range.  Thus
$W=J_PG^{-1/2}$ proves~\eqref{eq:g18-three-component-band}.

For completeness, the energy-kernel estimate uses no unrecorded second Riesz
mark.  In the gap-one coefficient normalization put
$\mathbb D_u=D+F_u$ and write $c_t=\mathcal P_u^-w_t-w_t$.  Since
$Dw_t=4w_t$ and $\mathbb D_u$ commutes with $\mathcal P_u^-$,
\[
 \mathbb D_u\mathcal P_u^-w_t=4w_t+k_t,
 \qquad k_t=Dc_t+F_uw_t+F_uc_t .
\]
The rooted relative bound and~\eqref{eq:g18-rooted-projection} give
\[
 \sum_J\eta^{-d(J,S_t)}\|(k_t)_J\|
 \le B_\eta(g):=A_\eta(g)+4a_\eta(g)
                  +a_\eta(g)A_\eta(g)\longrightarrow0 .
\]
The GNS intertwining~\eqref{eq:g18-cmp1-generator-intertwining} and
orthogonality of the Riesz projection give the one-mark identity
\[
 \bigl\langle\psi_s,\widehat H_\infty^-(u)\psi_t\bigr\rangle
 =\omega_u\!\left(\widehat w_s^*\,
        \widehat{\mathbb D_u\mathcal P_u^-w_t}\right),
 \qquad \widehat H_\infty^-(u)=\frac32H_\infty^-(u).
\]
Applying the same finite-contact/cluster-tail split to $4w_t+k_t$ now shows
that this energy kernel belongs to the same exponentially weighted
convolution algebra.  The convergent binomial series for $G^{-1/2}$ in that
algebra then puts $W^*H_\infty W$ there as well, proving analyticity of its
Bloch matrix.

For the literal source put
\[
 D_s=c_WO_s^W-\widehat w_s,\qquad
 M_D=3\sqrt2+1.
\]
Character orthogonality gives $c_W=i\sqrt2$ and
$D_s\Omega_0=0$.  Hence
\[
 \|\pi_u(D_s)\Omega_u\|^2
 \le6gM_D^2 .
\]
Define the three synthesis maps by
\[
 T_cf=\sum_s f_s\psi_s,
 \qquad
 T_Df=\sum_s f_sP_{F,\infty}^-(u)\pi_u(D_s)\Omega_u,
 \qquad T_W=T_c+T_D .
\]
Thus $T_W$ is exactly the projected, $c_W$-normalized literal Wilson source.
Its cross kernel is $C=T_c^*T_D$.  For contacting roots, the displayed
$O(\sqrt g)$ norm bound and Cauchy--Schwarz control $C_{st}$.  For separated
roots, use $\psi_s=\Gamma_u(w_s+c_s)$: charge oddness removes both one-point
terms, while the cluster bound controls $w_s$ and the rooted estimate controls
$c_s$.  The same contact/tail summation therefore gives
$\|C\|_{\mu,\sharp}\to0$.

Since $G=T_c^*T_c$ is invertible, put $E=G^{-1}C$.  Both $T_c$ and $T_D$ take
values in the target band, and $T_c$ is onto it.  Hence the exact identities
\[
 T_cG^{-1}T_c^*=I_{\operatorname{Ran}P_{F,\infty}^-},
 \qquad
 T_D=T_cG^{-1}T_c^*T_D=T_cE,
 \qquad
 T_W=T_c(I+E)
\]
show that $T_W$ is onto and bounded below once
$\|E\|_{\mu,\sharp}<1$.  Taking Gram matrices,
$G_W=(I+E)^*G(I+E)$, proves
~\eqref{eq:g18-wilson-frame}.
\end{proof}

\begin{remark}[Exact meaning of the closed result]
Theorem~\ref{thm:g18-fixed-spacing-band} closes the complete
three-component plaquette-shell band and the literal Wilson
\emph{source} residue at sufficiently small strong coupling and fixed
lattice spacing.  It does not isolate the one-dimensional $BB\dg$ kernel
inside each nonzero-momentum fiber, and it does not identify the Hamiltonian
semigroup with the isotropic Wilson-action transfer matrix.
\end{remark}

There is also a sharp scalar consequence for the compact carrier seed.  For an
oriented plaquette $p$, define the normalized self-adjoint charge-odd source
\[
 A_p=\frac{\chi_p-\bar\chi_p}{i\sqrt2}.
\]
Let $c_x$ be the elementary cube based at $x$, with outward incidence signs
$s(c_x,p)$, and put
\begin{equation}
 C_x=\frac1{\sqrt6}\sum_{p\subset\partial c_x}s(c_x,p)A_p,
 \qquad
 \eta_x(u)=P_{F,\infty}(u)\pi_u(C_x)\Omega_u .
 \label{eq:g18-cube-source}
\end{equation}
Its projected Gram kernel is
\begin{equation}
 G_u(x)=\langle\eta_0(u),\eta_x(u)\rangle .
 \label{eq:g18-gram-kernel}
\end{equation}
At $u=0$, plaquette orthogonality and
Proposition~\ref{prop:fixed-k-carrier} give the exact symbol
\begin{equation}
 \widehat G_0(k)=\sum_xe^{-ik\cdot x}G_0(x)=\frac{q(k)}6 .
 \label{eq:g18-free-gram}
\end{equation}

\begin{corollary}[Compact carrier-source residue]
\label{lem:g18-one-estimate}
Let $G_W(k)$ be the three-orientation literal-Wilson-source Gram symbol in
~\eqref{eq:g18-wilson-frame}.  Then the cube-boundary source
in~\eqref{eq:g18-cube-source} obeys
\begin{equation}
 \widehat G_u(k)
 \ge(1-g_0(u))
       (1-\|E\|_{\mu,\sharp})^2\frac{q(k)}6 .
 \label{eq:g18-positive-residue}
\end{equation}
In particular, on every $U\Subset\mathbb T^3\setminus\{0\}$ with
$q(k)\ge q_*>0$, the compact carrier source has residue at least
$(1-g_0(u))(1-\|E\|_{\mu,\sharp})^2q_*/6$, uniformly for $k\in U$.
\end{corollary}

\begin{proof}
In the three-orientation plaquette fiber, the normalized cube boundary has
coefficient vector $w(k)/\sqrt6$.  Theorem~\ref{thm:fact} gives
$\|w(k)\|^2=q(k)$, while the literal source in
~\eqref{eq:g18-cube-source} has the Wilson Gram, not the canonical-creator
Gram.  Equation~\eqref{eq:g18-wilson-frame} gives
$G_W(k)\succeq(1-g_0(u))(1-\|E\|_{\mu,\sharp})^2I_3$.  Therefore
\[
 \widehat G_u(k)
 =\frac16w(k)^*G_W(k)w(k)
 \ge(1-g_0(u))(1-\|E\|_{\mu,\sharp})^2\frac{q(k)}6 .
\]
\end{proof}

The full three-component band and its compact carrier-source residue are
therefore closed at fixed spacing.  The following continuation identifies the
second-order coefficient of the exact orthonormal symbol and closes its
rank-one internal sheet on every momentum ball compactly separated from
$\Gamma$.

% Publication-ready G18 continuation.  This file is input by
% G18_FIXED_SPACING_CARRIER_BRIDGE_INSERT.tex and assumes its repaired CMP(1)
% close-projection/GNS transport lemma and complete fixed-spacing
% three-component band theorem.
%
% Bibliography addition required if this insert is merged:
% \bibitem{Yarotsky2006GroundStates} D. A. Yarotsky,
%   \emph{Ground states in relatively bounded quantum perturbations of
%   classical lattice systems}, Commun. Math. Phys. \textbf{261}, 799--819
%   (2006), \href{https://arxiv.org/abs/math-ph/0412040}
%   {arXiv:math-ph/0412040}.

\subsection{Exact isolation of the internal carrier sheet away from
\texorpdfstring{$\Gamma$}{Gamma}}
\label{sec:g18-internal-sheet-closure}

This subsection works at fixed lattice spacing in the $SU(3)$ Hamiltonian
theory and in the unrescaled energy convention of Eq.~\eqref{eq:H}.  It uses
the repaired close-projection theorem and
Theorem~\ref{thm:g18-fixed-spacing-band}; no Wilson-action transfer matrix is
used.  The result closes the internal rank-one sheet on every fixed momentum
patch whose closure avoids $\Gamma$.  It does not give a gap uniform as
$k\to0$.

For $\mu>0$, let $\mathcal A_\mu$ denote the unital Banach $*$-algebra of
$3\times3$ convolution kernels $K=(K_{\alpha\beta}(x))$ on $\mathbb Z^3$ with
norm
\begin{equation}
 \|K\|_{\mu,\sharp}:=
 \max\left\{
 \max_\alpha\sum_{x,\beta}e^{\mu|x|_1}|K_{\alpha\beta}(x)|,
 \max_\beta\sum_{x,\alpha}e^{\mu|x|_1}|K_{\alpha\beta}(x)|
 \right\}.
 \label{eq:g18-weighted-algebra}
\end{equation}
Its Fourier transform is analytic for
$|\operatorname{Im}k_j|<\mu$ and
$\sup_{k\in\mathbb T^3}\|K(k)\|\le\|K\|_{\mu,\sharp}$.

\begin{lemma}[Coupling-analytic exponentially local symbol]
\label{lem:g18-u-analytic-symbol}
There are $\rho_0>0$ and $\mu>0$ such that the orthonormal-frame Hamiltonian
of Theorem~\ref{thm:g18-fixed-spacing-band} is represented by a map
\begin{equation}
 u\longmapsto h_u\in\mathcal A_\mu
 \label{eq:g18-banach-holomorphic}
\end{equation}
which is holomorphic for complex $|u|<\rho_0$.  For real $u$, $h_u(k)$ is
Hermitian.  The same statement holds for the canonical Gram kernel $G_u$ and
the physical one-mark energy kernel $\mathcal K_u$, and
\begin{equation}
 h_u=G_u^{-1/2}\mathcal K_uG_u^{-1/2}.
 \label{eq:g18-h-from-gk}
\end{equation}
Here the inverse square root is the branch satisfying $G_0^{-1/2}=I$.
\end{lemma}

\begin{proof}
Choose $\rho_0$ strictly inside both the repaired CMP(1) smallness domain and
the complex ground-state contraction domain.  With
$\lambda_0=27c_2\rho_0<1/33$, one may use the fixed contour
\[
 \gamma_{\rho_0}:\quad |z-4|=(1-\lambda_0)/8
\]
in the gap-one rescaling.  It encloses the target cluster for every
$|u|\le\rho_0$ and stays in the coefficient resolvent set.  Indeed, the
target disc has radius at most $4\lambda_0$, while the nearest other free
centres are $7/2$ and $17/4$.  The inequalities separating all three discs
from this contour reduce on the upper side to
$33\lambda_0<1$ (and on the lower side to the weaker
$\lambda_0<1/9$).  For complex $u$ this statement uses the coefficient
Neumann-resolvent bound, whose absolute majorant depends on $|u|$, rather
than a self-adjoint spectral-ordering argument.  This fixed contour, rather
than a contour depending on the evaluation point $u$, is what makes the
following holomorphic calculation literal.

The product-ground-state fixed-point map used in
\cite{Yarotsky2004QP} is analytic in the bounded interaction.  On a closed
complex $u$-disc contained in the contraction domain its contraction
constant and all coefficient majorants depend only on $|u|$.  The analytic
contraction theorem therefore makes its fixed point, the dressed coefficient
perturbation $F_u$, and the Neumann resolvent series holomorphic in their
weighted coefficient norms.  Contour integration on $\gamma_{\rho_0}$ makes
the coefficient Riesz projection $\mathcal P_u^-$ holomorphic in every
slightly weaker rooted norm.

We record why the spatial estimate remains valid on a complex disc.  In the
proof of \cite[Theorem~1]{Yarotsky2006GroundStates}, Lemmas~3--4 give
holomorphic polymer activities whose absolute values, uniformly on a closed
parameter disc, obey a majorant of the form
\[
 |a_u(\chi)|\le C_\rho\varepsilon_\rho^{|\operatorname{supp}\chi|},
 \qquad \varepsilon_\rho<1,
\]
after the coupling disc is made small enough.  The number of polymers of
size $n$ through a specified point grows at most exponentially in $n$.
Consequently the connected-cluster series is absolutely and locally
uniformly convergent for complex $u$, which is the argument used there to
prove weak-star analyticity.  With two marked local insertions, every
connected cluster joining supports $I$ and $J$ has size bounded below by a
constant times $\operatorname{dist}(I,J)$.  Summing the same activity
majorant therefore gives the complex-disc estimate
\begin{equation}
 \sup_{|u|\le\rho}
 |\omega_u^{\mathbb C}(A^*B)
   -\omega_u^{\mathbb C}(A^*)\omega_u^{\mathbb C}(B)|
 \le C_\rho^{|I|+|J|}\tau_\rho^{\operatorname{dist}(I,J)}
       \|A\|\|B\|,
 \qquad \tau_\rho<1.
 \label{eq:g18-complex-cluster-majorant}
\end{equation}
Here $\omega_u^{\mathbb C}$ denotes the holomorphic continuation; for real
$u$ it is the physical ground state.  This is an absolute cluster bound, not
an inference from real-$u$ positivity or from pointwise weak-star
analyticity.

Put $\mathbb D_u:=D+F_u$ for the gap-one coefficient Hamiltonian
and use the one-mark identities
\begin{align}
 (G_u)_{st}
 &=\omega_u\!\left(\widehat w_s^*\,\widehat{\mathcal P_u^-w_t}\right),
 \label{eq:g18-one-mark-g-analytic}\\
 (\widehat{\mathcal K}_u)_{st}
 &=\omega_u\!\left(\widehat w_s^*\,
       \widehat{\mathbb D_u\mathcal P_u^-w_t}\right),
 \qquad \mathcal K_u:=\frac23\widehat{\mathcal K}_u.
 \label{eq:g18-one-mark-k-analytic}
\end{align}
They agree with the Hilbert-space Gram and energy kernels for real $u$ but,
unlike a two-mark norm formula, also give their holomorphic continuation to
complex $u$.  To make the summation explicit, write
$\mathcal P_u^-w_t=w_t+c_t(u)$.  The rooted Neumann majorant gives, uniformly
for $|u|\le\rho<\rho_0$,
\[
 \sum_J\eta^{-d(J,S_t)}\|D_Jc_{t,J}(u)\|\le A_{\eta,\rho},
\]
and the same calculation applied to
\[
 \mathbb D_u\mathcal P_u^-w_t
 =4w_t+Dc_t+F_uw_t+F_uc_t
\]
gives a rooted majorant $A'_{\eta,\rho}$.  Both coefficient families are
holomorphic.  Exact-support projections commute with charge conjugation, so
$w_t$, $c_{t,J}$ and every displayed energy correction are charge odd.  The
analytically continued ground-state functional is charge even; hence both
one-point terms in
\eqref{eq:g18-complex-cluster-majorant} vanish identically, for complex as
well as real $u$.  For separated roots apply
\eqref{eq:g18-complex-cluster-majorant} term by term and use
$\operatorname{dist}(S_s,S_t)\le
\operatorname{dist}(S_s,J)+d(J,S_t)$, exactly as in the real CMP(3) Schur
sum.  Contacting roots form a fixed finite list and are bounded directly by
the absolutely convergent marked-cluster series.  Choose $\eta$ and then the
disc so that
$\theta_\rho:=\max\{C_\rho\eta,\tau_\rho\}<1$.  Since the number of roots at
distance $n$ grows quadratically, any
$0<\mu_0<-\log\theta_\rho$ gives constants $C,\mu_0>0$ such that
\[
 \sum_{x\in\mathbb Z^3}e^{\mu_0|x|_1}
 \sup_{|u|\le\rho}\bigl(
   |G_{\alpha\beta}(x;u)|
   +|\mathcal K_{\alpha\beta}(x;u)|
 \bigr)\le C,
 \qquad \rho<\rho_0.
\]
Finite spatial truncations are holomorphic, and this summable supremum
majorant makes them converge locally uniformly in $\mathcal A_\mu$ for any
$0<\mu<\mu_0$.  Thus $G_u$ and $\mathcal K_u$ are Banach-valued holomorphic
maps.

After shrinking $\rho_0$ if necessary,
$\|G_u-I\|_{\mu,\sharp}<1$.  The binomial series for $G_u^{-1/2}$ converges
locally uniformly in the unital Banach algebra $\mathcal A_\mu$ and is
holomorphic.  Equation~\eqref{eq:g18-h-from-gk} proves the claim.  The
coefficient calculation in the gap-one Hamiltonian
$\widehat H=(3/2)H$ would give the non-scalar coefficient
$(3/2)t_3=5/408$.  Multiplying the completed symbol by $2/3$ returns exactly
$t_3=5/612$ in the unrescaled normalization of Eq.~\eqref{eq:H}.
\end{proof}

\begin{lemma}[Finite-to-infinite second-order coefficient]
\label{lem:g18-finite-infinite-u2}
In the orthonormal frame of
Theorem~\ref{thm:g18-fixed-spacing-band},
\begin{equation}
 [u^0]h_u(k)=\frac83I_3,\qquad
 [u^1]h_u(k)=I_3,\qquad
 [u^2]h_u(k)=\sigma_2I_3+\frac5{612}B(k)B(k)^\dagger,
 \label{eq:g18-exact-infinite-u2}
\end{equation}
where $\sigma_2$ is a real scalar.  If the separate $SU(3)$ diagonal ledger
quoted in Section~\ref{sec:su3} is adopted, then $\sigma_2=11/306$; the
internal-sheet conclusion below does not use this value.
\end{lemma}

\begin{proof}
First work in a sufficiently large finite periodic box ($L\ge5$), in the
physical charge-odd trivial-flux shell, and use a fixed contour valid on a
complex $u$-disc.  Kato parallel transport gives a holomorphic intertwiner
from the free plaquette shell to the interacting Riesz island; on the real
axis it is unitary.  Differentiating its reduced operator at $u=0$ gives the ordinary
degenerate reduced-resolvent coefficient
\[
 P_-(-V)Q_-(E_{F,3}-H_E)^{-1}Q_-(-V)P_-
   -[u^2]E_{\rm vac,L}(u)I,
\]
with the conventional scalar terms absorbed into the diagonal ledger.  This
is precisely the $Q_-=1_--P_-$ prescription of Eq.~\eqref{eq:heff}, rather
than an eigenvalue-only comparison.  The retained-shell theorem,
the first-order charge/determinant census, and
Corollary~\ref{cor:assembly} therefore give
\[
 k_{L,0}=\frac83I,\qquad k_{L,1}=I,\qquad
 k_{L,2}=\sigma_{2,L}I+\frac5{612}BB^\dagger.
 \label{eq:g18-finite-kato-u2}
\]

The canonical Gram orthonormalization used to define $h_{L,u}$ is another
holomorphic intertwiner which equals the free plaquette frame at $u=0$ and
is unitary for real $u$.  The two identifications differ by a holomorphic
near-identity similarity $Z(u)$, unitary on the real axis.  If
$Z_0=I$ and
\[
 a(u)=a_0I+ua_1I+u^2a_2+\sum_{n\ge3}u^na_n,
\]
then direct expansion gives
\[
 [u^2]\,Z(u)^{-1}a(u)Z(u)=a_2;
\]
all possible frame terms are commutators with $a_0I$ or $a_1I$.  Thus the
full second-order matrix, not merely its eigenvalues, is unchanged by the
orthonormal-frame choice.  This is the point at which scalarity of both the
zeroth- and first-order coefficients is essential.

It remains to justify passage to infinite volume without confusing torus
convolution with convolution on $\mathbb Z^3$.  Let $\mathcal A_{\mu,L}$ be
the corresponding weighted convolution algebra on $\mathbb Z_L^3$, using
the torus distance, and let
\[
 (\Pi_LK)([x])=\sum_{n\in\mathbb Z^3}K(x+nL)
\]
be periodisation.  Then $\Pi_L:\mathcal A_\mu\to\mathcal A_{\mu,L}$ is a
contractive algebra homomorphism.  Let $G_{L,u}$ and
$\mathcal K_{L,u}$ be the periodic finite-volume one-mark kernels.  The
coefficient fixed-point and marked-cluster series converge term by term to
their infinite-volume series for every fixed centred displacement, locally
uniformly on a smaller complex $u$-disc.  The periodic ground-state limit is
the same state as the open-box limit because the small-perturbation ground
state is unique \cite{Yarotsky2005Uniqueness}.  This first identifies the
two limits for real $u$; local uniform holomorphy and the identity theorem
then identify their complex continuations on the common disc.

The proof of Lemma~\ref{lem:g18-u-analytic-symbol} gives the stronger bound,
uniformly in $L$,
\[
 \sum_{[x]}e^{\mu_0d_L([x],0)}
 \sup_{|u|\le\rho}
 \bigl(|G_{L,\alpha\beta}(x;u)|
       +|\mathcal K_{L,\alpha\beta}(x;u)|\bigr)\le C_\rho.
\]
Fix $0<\mu<\mu_0$.  Split the torus norm at
$d_L([x],0)=R$, with $2R<L$.  The finite part of
$G_{L,u}-\Pi_LG_u$ (and of the energy kernel) converges locally uniformly by
the preceding termwise limit.  Both tails, including every periodised image,
are at most $C_\rho e^{-(\mu_0-\mu)R}$.  Letting first $L\to\infty$ and then
$R\to\infty$ proves
\[
 \|G_{L,u}-\Pi_LG_u\|_{\mu,\sharp,L}\longrightarrow0,\qquad
 \|\mathcal K_{L,u}-\Pi_L\mathcal K_u\|_{\mu,\sharp,L}
 \longrightarrow0
\]
locally uniformly in $u$.  On a smaller disc all Gram kernels obey a common
$\|G_{L,u}-I\|_{\mu,\sharp,L}\le r_G<1$.  The binomial series is therefore
uniform in $L$, and the homomorphism property of $\Pi_L$ gives
\[
 \|h_{L,u}-\Pi_Lh_u\|_{\mu,\sharp,L}\longrightarrow0
\]
locally uniformly.  This also shows explicitly that wrap-around products in
the finite-torus inverse square root are accounted for by periodisation, not
silently identified with infinite convolution.  Cauchy's formula on any
circle inside that disc now gives
\begin{equation}
 [u^2]h_u(x) =
 \lim_{L\to\infty}[u^2]h_{L,u}(x)
 \label{eq:g18-cauchy-coefficient-limit}
\end{equation}
for each fixed displacement $x$.  This is stronger than convergence at a
few real couplings: local uniformity on a complex contour is what permits
interchanging the thermodynamic limit with Taylor-coefficient extraction.

At order two every connected process is supported on one face or on two
faces sharing one link.  Hence, once the roots and this finite process
neighbourhood fit without wrapping, the finite-volume coefficient is exactly
the infinite-lattice coefficient.  Lemma~\ref{lem:process-complete} excludes
all other off-diagonal supports, and
Theorem~\ref{thm:allrank} assigns the surviving oriented shared-link element
$t_3=5/612$.  Translation and cubic symmetry make the remaining diagonal a
scalar.  The coefficient is finite range, so the entrywise identity also
holds in every $\mathcal A_\mu$.  This proves
\eqref{eq:g18-exact-infinite-u2}.
\end{proof}

Define
\begin{equation}
 s_{\le2}(u):=\frac83+u+\sigma_2u^2,\qquad t_3:=\frac5{612}.
 \label{eq:g18-s2-definition}
\end{equation}

\begin{theorem}[Exact fixed-spacing carrier sheet on a nonzero-momentum patch]
\label{thm:g18-exact-fixed-k-carrier}
Assume the repaired CMP(1) lemma and
Theorem~\ref{thm:g18-fixed-spacing-band}.  Let $U$ be a relatively compact
open momentum ball whose closure lies in
$\mathbb T^3\setminus\{0\}$, and put
\[
 q_*:=\inf_{k\in U}q(k)>0.
\]
Then there is $u_U>0$ such that for every real
$0<|u|<u_U$ the exact infinite-volume target band over $U$ has a rank-one
analytic Riesz subbundle, descended from $\ker B(k)^\dagger$, separated from
the other two target branches by
\begin{equation}
 \operatorname{gap}_{\rm int}(k,u)
 \ge\frac12t_3u^2q_*,
 \qquad k\in U.
 \label{eq:g18-closed-internal-gap}
\end{equation}
Its direct integral is a reducing, translation-invariant subspace unitarily
equivalent to $L^2(U)$, on which translations and the Hamiltonian act as
\begin{equation}
 (T_xf)(k)=e^{ik\cdot x}f(k),
 \qquad
 (H_{\rm car}(u)f)(k)=E_{\rm car}(k,u)f(k).
 \label{eq:g18-carrier-representation-closed}
\end{equation}
The function $E_{\rm car}(k,u)$ is real analytic in $k$ on $U$.
\end{theorem}

\begin{proof}
Choose $0<r<\rho_0$ in Lemma~\ref{lem:g18-u-analytic-symbol} and set
\[
 M_r:=\sup_{|\zeta|=r}\|h_\zeta\|_{\mu,\sharp}<\infty,
 \qquad C_r:=\frac{2M_r}{r^3}.
\]
Cauchy's estimate in $\mathcal A_\mu$, followed by summing the Taylor tail,
gives for $|u|\le r/2$
\begin{equation}
 h_u(k)=s_{\le2}(u)I_3+t_3u^2B(k)B(k)^\dagger+R_u(k),
 \qquad
 \sup_{k\in\mathbb T^3}\|R_u(k)\|
 \le C_r|u|^3.
 \label{eq:g18-uniform-cubic-remainder}
\end{equation}
Indeed, the $n$th Banach-valued Taylor coefficient has norm at most
$M_rr^{-n}$, and
$\sum_{n\ge3}(|u|/r)^n\le2|u|^3/r^3$ on this disc.

One may therefore take, in addition to the smallness thresholds required by
CMP(1)--CMP(3),
\begin{equation}
 u_U:=\min\left\{\frac r2,
          \frac{t_3q_*}{4C_r}\right\}>0.
 \label{eq:g18-explicit-uU}
\end{equation}
For $0<|u|<u_U$,
\[
 \sup_{k\in U}\|R_u(k)\|
 \le\frac14t_3u^2q_*.
\]
The comparison matrix in
\eqref{eq:g18-uniform-cubic-remainder} has a simple carrier eigenvalue
$s_{\le2}(u)$ and a double eigenvalue
$s_{\le2}(u)+t_3u^2q(k)$.  Weyl's inequality separates the corresponding
exact spectral clusters by at least
\[
 t_3u^2q(k)-2\|R_u(k)\|
 \ge\frac12t_3u^2q_*.
\]
The small contour surrounding the lower cluster therefore defines a rank-one
Riesz projection analytic in $k$.  Put
$P_0(k)=v_0(k)v_0(k)^\dagger$ with
$v_0(k)=w(k)/\sqrt{q(k)}$; this is the limiting comparison projection chosen
by the $BB^\dagger$ coefficient, not a spectral projection of the triply
degenerate $u=0$ fiber.  The same Riesz estimate and
$\|R_u\|/(t_3u^2q_*)=O(|u|/q_*)$ show that, after decreasing $u_U$ if
necessary,
$\sup_{k\in U}\|P_{\rm car}(k,u)-P_0(k)\|<1/2$.  Hence
\[
 v_u(k):=\frac{P_{\rm car}(k,u)v_0(k)}
 {\|P_{\rm car}(k,u)v_0(k)\|}
\]
is a globally defined real-analytic unit section on the whole momentum ball.
Fourier direct integration gives
\eqref{eq:g18-carrier-representation-closed}; analyticity of the simple
eigenvalue follows from analyticity of the matrix symbol.
\end{proof}

\begin{corollary}[No dark carrier on a fixed momentum patch]
\label{cor:g18-no-dark-fixed-patch}
After possibly decreasing $u_U$, the rank-one sheet in
Theorem~\ref{thm:g18-exact-fixed-k-carrier} is seen uniformly by a dressed
cube-boundary Wilson source.  More precisely, put
\[
 v_0(k):=\frac{w(k)}{\sqrt{q(k)}},\qquad k\in U,
\]
and let $S_W(k,u)$ be the normalized literal-Wilson projected synthesis map
written in the orthonormal band frame, with the fixed CMP(4) normalization
chosen so that $S_W(k,0)=I_3$.  If $P_{\rm car}(k,u)$ is the exact rank-one
projection, then $u_U$ may be chosen so that
\begin{equation}
 \inf_{k\in U}
 \bigl\|P_{\rm car}(k,u)S_W(k,u)v_0(k)\bigr\|^2
 \ge\frac14,
 \qquad 0<|u|<u_U.
 \label{eq:g18-fixed-patch-carrier-overlap}
\end{equation}
The number $1/4$ is an absolute projected-norm bound for this normalized
coefficient vector, not a spectral fraction normalized by the full raw-source
covariance.  Consequently every $f\in C_c^\infty(U)$ gives a rapidly decaying
spatial smearing of the literal Wilson plaquettes, with cube-boundary
incidence, whose projection onto the exact carrier sheet is nonzero and
volume independent.
\end{corollary}

\begin{proof}
At $u=0$ the three target levels coincide, so there is no spectrally defined
rank-one carrier projection.  Instead let
$P_0(k)=v_0(k)v_0(k)^\dagger$ be the limiting comparison projection selected
by the $BB^\dagger$ coefficient.  Then
$P_0(k)S_W(k,0)v_0(k)=v_0(k)$.  The remainder estimate
\eqref{eq:g18-uniform-cubic-remainder} and the order-$u^2$ gap imply, by the
finite-dimensional Riesz-projection estimate, that
\[
 \sup_{k\in U}\|P_{\rm car}(k,u)-P_0(k)\|\longrightarrow0
 \quad (u\to0).
\]
CMP(3)--CMP(4) make this explicit.  In the notation
$W=T_cG^{-1/2}$ and $T_W=T_c(I+E)$,
\[
 S_W=W^*T_W=G^{1/2}(I+E),
\]
so $\|S_W-I\|_{\mu,\sharp}\to0$.  Therefore
\[
 \|P_{\rm car}S_Wv_0\|
 \ge1-\|P_{\rm car}-P_0\|-\|S_W-I\|.
\]
Since $\overline U$ is compact, the right-hand side is at least $1/2$ after
shrinking $u_U$.
Squaring gives \eqref{eq:g18-fixed-patch-carrier-overlap}.  Fourier transform
then turns $f\in C_c^\infty(U)$ into a rapidly decaying spatial smearing.
The identity
$f(k)v_0(k)=w(k)f(k)/\sqrt{q(k)}$ identifies its coefficients as the
cube-boundary incidence combination of the three literal Wilson plaquettes.
\end{proof}

\begin{remark}[What has and has not been closed]
\label{rem:g18-internal-sheet-firewall}
Theorem~\ref{thm:g18-exact-fixed-k-carrier} removes the two hypotheses in the
former conditional fixed-$k$ statement: it proves both the exact
orthonormal-symbol coefficient
$[u^2]h_u=\sigma_2I+(5/612)BB^\dagger$ and a remainder uniform in momentum
on a nonzero-momentum patch.  The threshold $u_U$ is
existential because the Yarotsky contraction majorants are existential, and
it necessarily deteriorates as $q_*\downarrow0$.  No rank-one separation at
$\Gamma$ is claimed; the cubic $T_1$ triplet meets there.

For every $a_t>0$, the exact Hamiltonian semigroup restricts to this subspace
as multiplication by $e^{-a_tE_{\rm car}(k,u)}$.  This does \emph{not}
identify $e^{-a_tH_\infty(u)}$ with the finite-temporal-spacing isotropic
Wilson-action transfer matrix.  The anisotropic-time/operator equivalence and
every continuum-limit statement remain separate open problems.
\end{remark}

